% !PRINTED 395
\weberpage{395}
We form the real syzygetic complex triple in which all
double tangents must occur:
\[
\begin{array}{rl}
1) & x\,y,\ p\,q \\
2) & x\,p,\ y\,q \\
3) & x\,q,\ y\,p.
\end{array}
\]
If a fifth real double tangent $r$ is present, we may assume
that it occurs in complex 1). But then this complex must
contain a third real pair $r\,s$, and either a fourth real pair,
in which case Theorem 3. already holds, or a conjugate
pair $\varrho\varrho'$.

If complex 2) does not contain solely real double tangents —
a case in which Theorem 3. would likewise hold — then an
imaginary pair $\xi\eta$ must occur in 2), and accordingly we thus consider the
following complexes:
\[
\begin{array}{rl}
1) & x\,y,\ p\,q,\ r\,s,\ \varrho\varrho' \\
2) & x\,p,\ y\,q,\ \xi\eta \\
4) & x\,r,\ y\,s.
\end{array}
\]
The two complexes 2), 4), however, are \gesperrt{azygetic}, since e.g. $r$ does not
occur in 2) [because it occurs in 1)]. Consequently 4) can also contain
only one of the two $\xi, \eta$, and consequently, since the complex
4) is real, $\xi, \eta$ cannot be conjugate. From this there results a
conjugate double pair for each of 2) and 4), and we have:
\[
\begin{array}{rl}
2) & x\,p,\ y\,q,\ \xi\eta,\ \xi'\eta' \\
4) & x\,r,\ y\,s,\ \xi\zeta,\ \xi'\zeta'.
\end{array}
\]
wherein $\zeta, \zeta'$ again denotes a pair of conjugate double tangents.

Now we form the real complex, syzygetic with 2) as well
as with 4),
\[
5)\quad \xi\xi',\ \eta\eta',\ \zeta\zeta',
\]
which can contain none of the elements $x, y, p, q, r, s$, and further the
two conjugate imaginary complexes
\[
\begin{array}{rl}
6) & x\xi,\ p\eta,\ r\zeta \\
7) & x\xi',\ p\eta',\ r\zeta',
\end{array}
\]
and 5), 6), 7) form an \gesperrt{azygetic triple}. $y, q, s$ do not occur in 6)
and 7), for otherwise they would have to be azygetic with two each
of the $x, p, r$, which by 2) and 4) is impossible.
